\chapter{Unique Selling Points, In Detail}

A USP is only a USP if it satisfies three tests: (i) a competing team is unlikely to have
it; (ii) it cannot be retrofitted in the last four hours; (iii) it survives a hostile
question. Each of the ten below is written against those tests, with an explicit build cost
so the team can cut the list if time runs short. \textbf{USPs 1--5 are non-negotiable.
6--8 are strongly recommended. 9--10 are bonus.}

\section{USP 1 --- The Zero-Circumvention Architecture}

\begin{uspbox}
\bhead{Claim.} APIx is the only submission that a Government of India ministry could
actually deploy, because it never circumvents a technical access control and it records the
legal basis of every single price quote in a tamper-evident ledger.
\end{uspbox}

\bhead{What it is.} A four-rung source ladder (licensed APIs $\to$ partner feeds $\to$
robots-permitted public endpoints $\to$ commercial datasets), a Politeness Governor that
every fetch passes through, and an append-only hash-chained Provenance Ledger exposed via a
public \code{/provenance} endpoint.

\bhead{Why others won't have it.} The PS text \emph{invites} teams to build a CAPTCHA
solver and an IP rotator. Most will. They will demo it proudly, and a MoSPI evaluator ---
whose ministry cannot lawfully deploy such a thing --- will mark it down. Reading the
contradiction in the PS and resolving it deliberately is a comprehension advantage, not a
technical one, which is exactly why it cannot be copied on the day.

\bhead{How to demo it.} Live: hit a host whose \code{robots.txt} disallows the fare path.
Show the request \emph{not being issued} and a \code{ROBOTS\_BLOCKED} ledger entry appearing
instead. Then call \code{/provenance/\{id\}} on a published index number and walk the chain
back to the source, the directive, and the rate limit that applied. Thirty seconds; nobody
else will have anything like it.

\bhead{Build cost.} $\approx$6 hours. Politeness Governor 3 h, ledger + hash chain 2 h,
provenance endpoint 1 h.

\bhead{The hostile question.} \emph{``So you have less data.''} Answer: yes, and here is the
coverage ledger quantifying exactly how much less, and the index standard errors that
implies. Then: \emph{``inadmissible data has a coverage of zero, because it cannot be
published.''}

\section{USP 2 --- CPI-2024-Native by Construction}

\begin{uspbox}
\bhead{Claim.} APIx is not ``compatible with'' the CPI --- it is built on the CPI 2024
series' own mathematics, classification and linking arithmetic. Integration is a
configuration change, not a migration project.
\end{uspbox}

\bhead{What it is.} Jevons elementary aggregates and Young/Modified Laspeyres higher-level
aggregation, matching MoSPI FAQ Q.20 and Q.21 exactly. COICOP 2018 sub-class 07.3.3 tagging
throughout. A linking-factor implementation validated against MoSPI's published General-level
factors (Rural 0.5222 / Urban 0.5320 / Combined 0.5267). A weight vector that is an
\emph{injected parameter}, so the Price Statistics Division supplies routes and weights and
the module consumes them --- which is a literal reading of the PS phrase \emph{``an
index-construction module based on PSD given routes and weights.''}

\bhead{Why others won't have it.} It requires reading the CPI 2024 FAQ and the Expert Group
Report. Most teams will use a mean, or a Laspeyres because it is the one they remember from
class, and will not know that MoSPI switched to COICOP 2018 in February 2026.

\bhead{How to demo it.} A single toggle on the dashboard: ``APIx own base'' $\leftrightarrow$
``rebased to CPI 2024$=$100 via linking factor'', with the official air fare item index
overlaid. Say out loud: \emph{``Jevons at the bottom, Young at the top --- because that is
what you do.''}

\bhead{Build cost.} $\approx$4 hours, mostly in the index engine, which you are building
anyway.

\bhead{The hostile question.} \emph{``Why Jevons and not Dutot?''} Answer with the AM--GM
worked example from \S3.2, the time-reversal property, and the lognormality of fares. Have
the three-formula slide ready.

\section{USP 3 --- The Dual-Time-Basis Index}

\begin{uspbox}
\bhead{Claim.} APIx publishes both the \emph{acquisition-basis} index (what it costs to buy
a flight today) and the \emph{use-basis} index (the price of travel consumed this month),
and reconciles them. No Indian price series distinguishes the two.
\end{uspbox}

\bhead{What it is.} \S\ref{sec:dual}. The booking-basis series attributes a quote to the
collection date; the travel-basis series attributes it to the departure date. The UK ONS
records air fare changes \emph{in the month the flight departs, not when the ticket is
bought} --- so the travel basis is the internationally CPI-consistent one, while the booking
basis is what a consumer, a regulator or an RBI nowcaster actually wants.

\bhead{Why others won't have it.} Almost nobody knows the distinction exists. It comes
directly from reading foreign NSI methodology, and it is the clearest possible signal that
the team engaged with price-statistics literature rather than with scraping tutorials.

\bhead{How to demo it.} One chart, two lines, a basis toggle, and one sentence: \emph{``These
are the same quotes. This one answers `what does a ticket cost today'; this one answers
`what did air travel cost this month'. CPI needs the second. RBI wants the first. We
publish both.''}

\bhead{Build cost.} $\approx$2 hours. It is a re-attribution of the same fact table --- a
different \code{GROUP BY}, essentially free once the schema stores both $t$ and $d$.

\bhead{The hostile question.} \emph{``Which one is the headline?''} Answer: travel basis is
the headline because it is CPI-consistent; booking basis is published alongside and is the
one used for the nowcast.

\section{USP 4 --- The Availability-Adjusted Index}

\begin{uspbox}
\bhead{Claim.} When cheap fare buckets sell out, a naive matched index reports \emph{zero
inflation} during exactly the surge episodes consumers complain about. APIx detects this,
quantifies it, and corrects it with a continuous, monotone adjustment that reduces to
standard practice in the normal case.
\end{uspbox}

\bhead{What it is.} \S\ref{sec:avail}. The availability ratio $A_g^{(t)}$, a
lowest-available-fare index, and the blend
$I^{\text{adj}} = (I^{\text{matched}})^{A}(I^{\text{LAF}})^{1-A}$.

\bhead{Why others won't have it.} It requires noticing that missingness in fare data is
\emph{informative} --- a selection-on-the-dependent-variable problem. Teams that treat a
sold-out flight as a failed scrape will silently build the bias in and never see it.

\bhead{How to demo it.} The strongest ninety seconds in the deck. Pick a real surge period.
Three lines: matched (flat, wrong), LAF (spiking), adjusted (spiking, correct). Underneath,
$A_g^{(t)}$ as an area chart collapsing toward zero. Then: \emph{``This is the Kumbh problem
and the festival-fare problem. A naive index cannot see it. Ours can.''}

\bhead{Build cost.} $\approx$5 hours: availability capture in the parser (2 h), LAF index
(1 h), blend and charting (2 h).

\bhead{The hostile question.} \emph{``Is that blend principled or invented?''} Answer
honestly: it is a proposed estimator, not an established standard. Defend it on three
properties --- it is continuous in $A$, monotone, and exactly equals the matched index when
$A=1$, so it never disturbs the normal case. Then note that the alternative (dropping
sold-out cells) is \emph{also} not an established standard, and is provably biased. Offer
the LAF series as a published diagnostic either way. \textbf{Do not claim it is official
methodology.}

\section{USP 5 --- The Lead-Time Elasticity Surface}

\begin{uspbox}
\bhead{Claim.} APIx treats advance-purchase window as a first-class published dimension,
estimates the fare--lead-time elasticity $\eta(\tau)$ with standard errors, and reports the
optimal booking lead time $\tau^\star$ per route.
\end{uspbox}

\bhead{What it is.} The pooled fixed-effects log-linear model of \S4.9, the elasticity
$\eta(\tau)=\beta_1+2\beta_2\ln\tau$, the trough
$\tau^\star=\exp(-\beta_1/2\beta_2)$, and the Gaussian-mixture recovery of the underlying
fare ladder.

\bhead{Why others won't have it.} The PS asks for ``lead-time elasticity curves'' and most
teams will plot average fare against $\tau$ and call it an elasticity curve. An average is
not an elasticity. Running an actual regression with route/carrier/day-of-week/month fixed
effects and clustered standard errors takes half an hour and is a completely different
object.

\bhead{How to demo it.} The 3-D surface is the visual, but the \emph{number} is the payload:
\emph{``For \rt{del}--\rt{bom}, $\tau^\star$ is 26 days --- book about four weeks out.
Elasticity at one week is $-0.34$: a 10\% longer lead time buys you 3.4\% off.''} That is a
press headline, and press headlines are what MoCA and DGCA want from a transparency product.

\bhead{Build cost.} $\approx$4 hours. \code{statsmodels} regression 1 h, surface plot 2 h,
mixture model 1 h.

\bhead{The hostile question.} \emph{``Is that causal?''} Answer: no --- it is a conditional
association within a revenue-management system, and the fixed effects absorb route, carrier,
day-of-week and seasonal confounders but not unobserved demand. Say ``association'', never
``causal effect''. Knowing the difference is itself the answer.

\section{USP 6 --- The Rigorous Back-Test Harness}
\label{sec:backtest}

\begin{uspbox}
\bhead{Claim.} The PS requires 30 days of back-tested results against DGCA fare data. APIx
delivers that as a \emph{statistical agreement analysis}, not an eyeballed overlay.
\end{uspbox}

\bhead{What it is.} A harness that compares the APIx monthly movement against (a) DGCA
sector-average fare movements and (b) the CPI 2024 air fare item index, and reports:

\begin{center}
\footnotesize
\begin{tabularx}{\textwidth}{@{}L{4.4cm} X@{}}
\toprule
\bhead{Statistic} & \bhead{What it answers} \\
\midrule
Pearson and Spearman $\rho$ & Do they move together, in level and in rank? \\
Mean absolute deviation of the y-o-y rate & By how much do we differ, in inflation points? \\
Bland--Altman plot & Is the disagreement constant, or does it grow with the level? (The right tool for method-agreement, and almost nobody uses it) \\
Directional agreement rate & In what \% of months do both move the same way? Often the most persuasive single number \\
Theil's $U$ vs.\ random walk & Does APIx beat a naive forecast? \\
Diebold--Mariano test & Is the forecast improvement statistically significant? \\
\bottomrule
\end{tabularx}
\end{center}

\bhead{Why others won't have it.} Most will show two lines and say ``as you can see, they
track closely.'' A Bland--Altman plot and a Diebold--Mariano statistic is a different
register of seriousness.

\bhead{How to demo it.} Show the overlay, then say: \emph{``and here is the agreement
analysis''} and put up the statistics table. Report an unflattering number if you have one;
it buys enormous credibility.

\bhead{Build cost.} $\approx$4 hours.

\begin{warnbox}
\bhead{The 30-day problem, and how to be honest about it.} If collection starts at the
hackathon, there is no 30-day history. Run \textbf{two tracks} and label them unambiguously:
\begin{itemize}
\item \textbf{Track A (live).} Start collecting the moment the basket is defined --- ideally
weeks before the finale. Even 14 real days is worth more than 90 synthetic ones.
\item \textbf{Track B (replay).} A synthetic generator calibrated to published DGCA monthly
average-fare movements and to known seasonality, used to exercise the pipeline at scale.
\end{itemize}
Put a visible \textbf{SYNTHETIC} watermark on every Track B chart. If a judge discovers an
unlabelled synthetic chart, the entire submission's credibility is gone. If you label it,
you look disciplined.
\end{warnbox}

\section{USP 7 --- The CPI Transport Nowcast}

\begin{uspbox}
\bhead{Claim.} APIx is available daily; the CPI air fare item arrives monthly around the
12th. APIx therefore \emph{leads} the official series, and we quantify by how much.
\end{uspbox}

\bhead{What it is.} The mixed-frequency bridge regression of \S4.10, with an ATF control from
PPAC, evaluated out-of-sample against a random-walk benchmark, plus Granger-causality and
lead--lag cross-correlation.

\bhead{Why others won't have it.} It requires connecting the airfare work to the actual
monetary-policy use case. This is the USP that speaks to RBI rather than to MoSPI, and it
turns ``a dashboard'' into ``an economic indicator''.

\bhead{How to demo it.} A chart with the official monthly series as step points and the APIx
daily series continuous, with the nowcast for the not-yet-published month highlighted and
its confidence band shown.

\bhead{Build cost.} $\approx$3 hours --- but it is \textbf{blocked on having history}. If
Track A is short, present the methodology and the code, and label the result as
illustrative. Do not fabricate an RMSE.

\bhead{The hostile question.} \emph{``Air fare is a tiny weight. Who cares?''} Answer: agreed
on the weight --- Transport is 8.796\% of combined CPI and air fare is a slice of that. The
claim is not that this moves headline inflation. The claim is that the \emph{method}
generalises: the same pipeline nowcasts hotel tariffs, cab fares, and the twelve online
markets MoSPI just introduced. Air fare is the pilot, not the prize.

\section{USP 8 --- Hash-Chained Reproducibility}

\begin{uspbox}
\bhead{Claim.} Every published APIx number can be regenerated from raw quotes by a single
command, and the collection history is tamper-evident.
\end{uspbox}

\bhead{What it is.} \code{method\_version} and \code{weights\_version} stamped on every index
value; a \code{make reproduce} target that recomputes the entire published series from the
Parquet cold tier and diffs against \code{fact\_index\_value}; the hash-chained provenance
ledger.

\bhead{Why others won't have it.} Reproducibility is invisible until someone asks for it,
so nobody builds it under time pressure. But it is a first-order requirement for official
statistics --- an index number that cannot be audited cannot be published by a ministry.

\bhead{How to demo it.} Run \code{make reproduce} live. It prints \code{0 differences across
N index values}. Then flip one byte in a stored quote and re-run: the hash chain breaks and
the harness reports exactly which record and which downstream index values are affected.

\bhead{Build cost.} $\approx$3 hours if the schema was designed for it from the start;
essentially impossible to bolt on afterwards. \textbf{Decide this in hour one.}

\section{USP 9 --- The Frugal Deployment Envelope}

\begin{uspbox}
\bhead{Claim.} APIx runs on a single modest VM with zero commercial proxy spend, and is
deployable on NIC/MeghRaj-class government infrastructure.
\end{uspbox}

\bhead{What it is.} The costing in \S6.9. Roughly 1.8 MB of raw quotes per day, under a GB
per year compressed, one small VM for collection and one for the database.

\bhead{Why it matters.} A scraping-first architecture needs a residential proxy pool, which
costs real money \emph{per gigabyte} and is a recurring line item a ministry must justify. It
is also the exact expenditure that makes a compliance reviewer uncomfortable. The
compliance choice and the cost choice are the same choice --- point that out; it is a nice
piece of narrative economy.

\bhead{How to demo it.} A single slide: monthly cost, storage growth curve, and the
statement \emph{``no external paid dependency in the critical path.''}

\bhead{Build cost.} $\approx$1 hour to instrument and chart.

\section{USP 10 --- The Fare Transparency Layer}

\begin{uspbox}
\bhead{Claim.} APIx implements, as a working prototype, two of the three recommendations the
Parliamentary Standing Committee on Transport, Tourism and Culture made to the Ministry of
Civil Aviation in its 375th Report.
\end{uspbox}

\bhead{What it is.} The Committee, reporting on Demands for Grants 2025--26, found the DGCA
\emph{lacks capacity to proactively monitor airfare trends} and recommended (i) an AI-driven
predictive fare monitoring system, (ii) a public \textbf{``Pricing Transparency Index''}
rating airlines on adherence to fair pricing, and (iii) an \textbf{``Airfare Vigil''}
consumer reporting app. APIx's route-level index, elasticity surface and fare-ladder
recovery are (i); a per-carrier tariff-band compliance score against Rule 135 of the Aircraft
Rules, 1937 is (ii).

\bhead{Why others won't have it.} It requires knowing the report exists. This is a research
advantage, and it lets the pitch end on \emph{``this is not a hackathon idea; this is a
standing Parliamentary recommendation, prototyped.''}

\bhead{How to demo it.} A consumer-facing panel: for each carrier and route, the observed
fare distribution against the declared tariff band, an outlier count, and a transparency
score. Frame it as \emph{additional} to the CPI mandate, not a substitute --- MoSPI owns the
PS, and the CPI work is the main line.

\bhead{Build cost.} $\approx$3 hours, and it depends on obtaining declared tariff bands, which
may not be publicly available. \textbf{If the bands are not obtainable, build the panel on
observed dispersion instead and say so.}

\section{USP summary: build order under time pressure}

\begin{center}
\footnotesize
\begin{tabularx}{\textwidth}{@{}C{0.7cm} L{4.6cm} C{1.5cm} C{1.7cm} X@{}}
\toprule
\bhead{\#} & \bhead{USP} & \bhead{Hours} & \bhead{Priority} & \bhead{Blocked on} \\
\midrule
1 & Zero-Circumvention Architecture & 6 & \textbf{Must} & Nothing \\
2 & CPI-2024-Native & 4 & \textbf{Must} & Weight data \\
3 & Dual-Time-Basis & 2 & \textbf{Must} & Schema decision \\
4 & Availability-Adjusted Index & 5 & \textbf{Must} & Availability capture in parser \\
5 & Lead-Time Elasticity Surface & 4 & \textbf{Must} & $\ge$2 weeks of data \\
8 & Hash-Chained Reproducibility & 3 & High & \textbf{Hour-one schema decision} \\
6 & Back-Test Harness & 4 & High & DGCA history \\
9 & Frugal Deployment Envelope & 1 & High & Instrumentation \\
7 & CPI Transport Nowcast & 3 & Medium & History length \\
10 & Fare Transparency Layer & 3 & Bonus & Tariff bands \\
\bottomrule
\end{tabularx}
\end{center}

\noindent Total for the five ``Must'' USPs: \textbf{21 hours}, on top of the base build. That
is the realistic scope. Cutting 7 and 10 costs nothing in the pitch; cutting 1 or 4 costs the
whole differentiation.
